% =============================================================================
% Vehicle Parameter Identification — Report
% =============================================================================
\documentclass[11pt,a4paper]{article}
\usepackage{amsmath,amssymb}
\usepackage{siunitx}
\usepackage{booktabs}
\usepackage{graphicx}
\graphicspath{{../figures/}}
\usepackage[margin=2.5cm]{geometry}
\usepackage{hyperref}
\usepackage{float}

% Conditionally include a figure only if the PDF exists
\newcommand{\includefigure}[3][width=0.7\textwidth]{%
  \IfFileExists{../figures/#2}{%
    \begin{figure}[H]
      \centering
      \includegraphics[#1]{#2}
      \caption{#3}
    \end{figure}%
  }{\par\textbf{[Figure not yet generated: \detokenize{#2}]}\par}%
}

\title{Vehicle Parameter Identification Report}
\author{}
\date{}

\begin{document}
\maketitle

\section{Vehicle Parameter Identification}\label{sec:param-id}

The dynamic bicycle model used by the MPC controller requires physical
parameters that cannot be taken from datasheets alone.
This section describes the experimental procedures used to identify
them on the real F1/10th vehicle (Traxxas Slash~4$\times$4 chassis,
VESC~6 controller, Kv\,=\,3500 motor).

All tests use a common ROS\,2 framework that records time-stamped
sensor data at approximately \SI{200}{\hertz}:
VESC odometry (speed from ERPM), the onboard IMU (3-axis accelerometer
and gyroscope), and motor state (current, voltage).
Each test is run five times; the file with the highest sample count
is used for reporting.  Raw data is saved to CSV for post-processing.

\paragraph{Vehicle constants (measured / CAD).}
\begin{table}[H]
\centering
\begin{tabular}{lll}
\toprule
\textbf{Quantity} & \textbf{Value} & \textbf{Source} \\
\midrule
Total mass $m$           & \SI{3.314}{\kilogram}  & Measured (scale) \\
Wheelbase $L$            & \SI{0.324}{\metre}     & CAD model \\
CG to front axle $l_f$  & \SI{0.166}{\metre}     & CAD model \\
CG to rear axle $l_r$   & \SI{0.160}{\metre}     & CAD model \\
Effective wheel radius $r_\text{eff}$ & \SI{0.051}{\metre} & Measured \\
Gear ratio $G$           & $11.82$                & Datasheet \\
Motor $K_v$              & $3500\;\text{RPM/V}$   & Datasheet \\
\bottomrule
\end{tabular}
\end{table}

%%% ===================================================================
\subsection{Dynamic Bicycle Model}\label{sec:bicycle-model}

The MPC uses a 7-state dynamic bicycle model with state
$\mathbf{x} = [x,\, y,\, \psi,\, v_x,\, v_y,\, \omega,\, \omega_w]$
and control inputs $\mathbf{u} = [\delta,\, T_\text{motor}]$.
The equations of motion are:
\begin{align}
  \dot{v}_x &= \frac{1}{m}\bigl(F_x - F_{yf}\sin\delta + m\,v_y\,\omega\bigr)
  \label{eq:dvx}\\
  \dot{v}_y &= \frac{1}{m}\bigl(F_{yf}\cos\delta + F_{yr} - m\,v_x\,\omega\bigr)
  \label{eq:dvy}\\
  \dot{\omega} &= \frac{1}{I_z}\bigl(l_f\,F_{yf}\cos\delta - l_r\,F_{yr}\bigr)
  \label{eq:domega}\\
  \dot{\omega}_w &= \frac{1}{I_w}\Bigl(\frac{T_\text{motor}}{G} - F_x\,R_w\Bigr)
  \label{eq:domegaw}
\end{align}
where $G$ is the gear ratio, $R_w$ the wheel radius, and $I_w$ the
drivetrain inertia. Lateral tire forces use a linear model:
\begin{equation}\label{eq:tire-lat}
  F_{yf} = C_{\alpha f}\,\alpha_f, \qquad
  F_{yr} = C_{\alpha r}\,\alpha_r
\end{equation}
with slip angles
\begin{equation}\label{eq:slip-angles}
  \alpha_f = \delta - \arctan\!\frac{v_y + l_f\,\omega}{v_x}, \qquad
  \alpha_r = -\arctan\!\frac{v_y - l_r\,\omega}{v_x}.
\end{equation}
The longitudinal tire force is $F_x = C_x\,\kappa$, where $\kappa$ is the
slip ratio. Table~\ref{tab:params} lists all model parameters and the
test used to identify each one.

%%% ===================================================================
\subsection{Tire Friction Coefficient}\label{sec:friction}

\paragraph{Theory.}
At steady-state circular driving the centripetal acceleration is
$a_y = v^2/R$. The maximum sustainable $a_y$ before the tires begin to
slide gives the friction coefficient:
\begin{equation}\label{eq:mu}
  \mu = \frac{a_{y,\max}}{g}.
\end{equation}
The IMU measures $a_y$ directly, independent of odometry errors.

\paragraph{Procedure.}
The car drives circles at a fixed steering angle
($\delta_\text{cmd} = \SI{0.3}{\radian}$) with speed increasing
in \SI{0.5}{\metre\per\second} steps from \SIrange{2.0}{5.0}{\metre\per\second}.
At each speed, \SI{8}{\second} of steady-state IMU data is recorded.

\paragraph{Data.}
\includefigure{friction_ay_vs_speed.pdf}{%
  Lateral acceleration (IMU) vs.\ speed, grouped by speed step.
  The friction coefficient $\mu$ is determined from the
  saturation plateau at high speed.\label{fig:friction}}

\begin{table}[H]
\centering
\caption{Friction test results: mean lateral acceleration and friction
  coefficient at each speed step (5 runs, $\sim$32\,000 samples per
  speed).}\label{tab:friction}
\begin{tabular}{S[table-format=1.1]
                S[table-format=1.2] @{${}\pm{}$} S[table-format=1.2]
                S[table-format=1.3]}
\toprule
{Speed (m/s)} & \multicolumn{2}{c}{$\bar{a}_y$ (m/s$^2$)} & {$\mu = a_y/g$} \\
\midrule
2.0  & 3.92 & 0.76 & 0.400 \\
2.5  & 5.39 & 1.06 & 0.549 \\
3.0  & 6.43 & 1.22 & 0.655 \\
3.5  & 6.89 & 1.46 & 0.702 \\
4.0  & 7.25 & 1.59 & 0.739 \\
4.5  & 7.55 & 1.89 & 0.769 \\
5.0  & 7.72 & 2.46 & 0.787 \\
\bottomrule
\end{tabular}
\end{table}

The measured $\mu$ increases with speed and begins to plateau at
$v \geq \SI{4.5}{\metre\per\second}$, approaching $\mu \approx 0.79$.
This value is consistent with rubber tires on smooth indoor flooring.
The increasing standard deviation at higher speeds reflects the onset
of intermittent sliding as the traction limit is approached.

\paragraph{Validation.}
The friction coefficient is cross-checked against the motor torque
test (Section~\ref{sec:motor-torque}), where the traction force
saturates at $F_\text{sat} \approx \SI{24}{\newton}$, giving
$\mu_x = F_\text{sat}/(m\,g) = 24/(3.314 \times 9.81) \approx 0.74$.
Both estimates are consistent, confirming $\mu \approx 0.74$--$0.79$
depending on the loading direction (longitudinal vs.\ lateral).

\paragraph{Uncertainty.}
\begin{itemize}
  \item IMU bias: \SI{2}{\second} stationary calibration; residual
        $< \SI{0.05}{\metre\per\second\squared}$.
  \item Result is surface-specific (smooth indoor floor); $\mu$ will
        differ on carpet, concrete, or asphalt.
  \item Tire temperature affects grip; all tests run at ambient temperature.
\end{itemize}

%%% ===================================================================
\subsection{Cornering Stiffness}\label{sec:cornering}

\paragraph{Theory.}
In steady-state cornering $(\dot{v}_y=\dot{\omega}=0)$, the front and
rear lateral forces and slip angles reduce to:
\begin{equation}\label{eq:Fy-ss}
  F_{yf} = m\,a_y\,\frac{l_r}{L}, \qquad
  F_{yr} = m\,a_y\,\frac{l_f}{L},
\end{equation}
\begin{equation}\label{eq:alpha-ss}
  \alpha_f \approx \delta - \frac{l_f\,\omega}{v_x}, \qquad
  \alpha_r \approx \frac{l_r\,\omega}{v_x},
\end{equation}
where $\omega$ is the IMU yaw rate and $v_x$ the ERPM-based speed.
The cornering stiffnesses follow from
$C_{\alpha f} = F_{yf}/\alpha_f$ and $C_{\alpha r} = F_{yr}/\alpha_r$.

\paragraph{Procedure.}
Same steady-state circles as the friction test, at speeds from
\SIrange{1.0}{3.5}{\metre\per\second}. At each speed, \SI{5}{\second} of
$v_x$, $\omega$ (IMU), and $a_y$ (IMU) are recorded.

\paragraph{Data.}
\includefigure[width=0.48\textwidth]{cornering_Fy_vs_alpha.pdf}{%
  Lateral force vs.\ slip angle for front and rear axles.  The linear
  region yields the cornering stiffness $C_\alpha$.\label{fig:cornering-Fy}}

\includefigure[width=0.48\textwidth]{cornering_alpha_vs_speed.pdf}{%
  Slip angle growth with speed for front and rear axles.\label{fig:cornering-alpha}}

\paragraph{Limitation.}
The current data was collected at a single steering angle
($\delta_\text{cmd} = \SI{0.1}{\radian}$), which produces a narrow
range of slip angles.  This limits the confidence in the fitted
cornering stiffness, particularly for the front axle where the
front slip angle range is very small.
The test script has been modified to sweep multiple steering angles
(0.08, 0.12, 0.16, 0.20, \SI{0.24}{\radian}) in future runs
to improve identification.

\paragraph{Uncertainty.}
\begin{itemize}
  \item Relies on measured $m$, $l_f$, $l_r$
        ($\pm 1$--$\SI{2}{\milli\metre}$).
  \item Single steering angle limits the observable slip-angle range.
  \item Servo gain mismatch: the actual steering angle is approximately
        86\% of the commanded angle (see Section~\ref{sec:wheelbase}),
        which shifts the absolute $C_\alpha$ estimate.
  \item ERPM noise affects $v_x$ and thus slip-angle calculation.
\end{itemize}

%%% ===================================================================
\subsection{Longitudinal Tire Stiffness}\label{sec:long-stiffness}

\paragraph{Theory.}
The slip ratio is
\begin{equation}\label{eq:kappa}
  \kappa = \frac{v_\text{wheel} - v_\text{body}}
               {\max(|v_\text{wheel}|,\,|v_\text{body}|)},
\end{equation}
where $v_\text{wheel}$ comes from ERPM and $v_\text{body}$ from an
independent velocity sensor. In the linear region ($|\kappa|<0.15$):
\begin{equation}\label{eq:Fx}
  F_x = C_x\,\kappa, \qquad F_x = m\,a_x.
\end{equation}

\paragraph{Limitation.}
Because the only available body-velocity estimate is obtained by
integrating the VESC IMU accelerometer, and the reading is
affected by chassis pitch during braking (gravity
projects onto the longitudinal axis), $v_\text{body}$ drifts rapidly.
The computed slip ratio $\kappa$ ranges from
$-1.0$ to $+0.9$ --- physically impossible values that confirm the
integration error.
$C_x$ cannot be reliably identified from the current data;
a lidar-based scan-matching odometry would provide an independent
body velocity to resolve this issue.

\paragraph{Data.}
\includefigure[width=0.48\textwidth]{long_v_comparison.pdf}{%
  $v_\text{wheel}$ (ERPM) and $v_\text{body}$ (IMU-integrated) vs.\ time.
  The divergence during braking is caused by IMU pitch
  contamination.\label{fig:long-v}}

\includefigure[width=0.48\textwidth]{long_Fx_vs_kappa.pdf}{%
  Traction force balance from the motor torque test: motor-side force
  $F_\text{motor}$ vs.\ body-side force $F_\text{body} = m\,a_x$.
  The linear slope gives drivetrain efficiency; the saturation
  plateau indicates the traction limit $\mu_x$.\label{fig:long-Fx}}

\paragraph{Traction force balance.}
Although $C_x$ could not be identified, the motor torque test data
(Section~\ref{sec:motor-torque}) provides a traction force balance.
The wheel-side force is computed from motor current as
$F_\text{motor} = K_{t,\text{eff}} \cdot I$, where
$K_{t,\text{eff}} = \SI{0.583}{\newton\per\ampere}$ is the measured
effective current-to-force gain.  The body-side force is
$F_\text{body} = m\,a_x$ from the IMU.  In the linear (unsaturated)
regime, the two forces are proportional with an efficiency of
$\eta \approx 0.92$ compared to the theoretical
$K_t\,G/r_\text{eff} = \SI{0.632}{\newton\per\ampere}$.  At high
current the body acceleration saturates, revealing the longitudinal
traction limit $\mu_x \approx 0.74$.

%%% ===================================================================
\subsection{Motor Torque Mapping}\label{sec:motor-torque}

\paragraph{Theory.}
Two independent estimates of wheel force:
\begin{align}
  F_\text{current} &= \frac{K_t \cdot G}{r_\text{eff}}\,I_\text{motor},
  \label{eq:F-current} \\
  F_\text{accel}   &= m\,a_x \quad\text{(IMU)},
  \label{eq:F-accel}
\end{align}
where $K_t = 60/(2\pi\,K_v) = \SI{0.00273}{\newton\metre\per\ampere}$
and $G=11.82$.  Plotting $F_\text{accel}$ against $I_\text{motor}$
yields the effective current-to-force gain (slope) and allows
identification of the VESC current limit.

\paragraph{Procedure.}
Accelerate from rest to target speeds
(\num{1.5}, \num{2.0}, \num{2.5}, \SI{3.0}{\metre\per\second}),
recording motor current and IMU $a_x$ simultaneously.

\paragraph{Data.}
\includefigure[width=\textwidth]{torque_I_vs_F.pdf}{%
  Motor torque mapping.
  \emph{Left:} raw scatter of all acceleration-phase data,
  coloured by velocity.  The red shaded region marks the
  VESC current limit ($\sim$\SI{65}{\ampere}).
  \emph{Right:} bin-averaged force vs.\ current in the linear region
  (5--\SI{55}{\ampere}), with the measured fit and theoretical
  line.\label{fig:torque-IF}}

\includefigure[width=0.7\textwidth]{torque_current_vs_time.pdf}{%
  Motor current (top) and commanded speed (bottom) vs.\ time for
  each speed step.  Green/red shading indicates
  acceleration/braking phases.  Time gaps between steps are
  collapsed.\label{fig:torque-time}}

\begin{table}[H]
\centering
\caption{Motor torque test: per-speed current and force statistics
  during the acceleration phase.}\label{tab:torque}
\begin{tabular}{S[table-format=1.1]
                S[table-format=2.1]
                S[table-format=2.1]
                S[table-format=1.1]
                S[table-format=4.0]}
\toprule
{Cmd speed (m/s)} & {$\bar{I}$ (A)} & {Peak $I$ (A)} & {$\bar{F}_\text{net}$ (N)} & {Samples} \\
\midrule
1.5 & 16.9 & 70.0 & 1.5  & 3204 \\
2.0 & 18.7 & 70.6 & 2.0  & 3236 \\
2.5 & 19.6 & 69.0 & 2.4  & 3235 \\
3.0 & 21.5 & 69.7 & 2.8  & 3233 \\
\bottomrule
\end{tabular}
\end{table}

The mean current and force increase with speed as expected.
The peak current of $\sim$\SI{70}{\ampere} at every speed step
corresponds to the VESC current limit during the initial launch
from standstill.  The mean force is low because the average
includes the steady-state cruising period where the car has reached
its target speed and acceleration is near zero.

\paragraph{Key results.}
\begin{itemize}
  \item \textbf{Measured $K_{t,\text{eff}}$}:
        $\SI{0.583}{\newton\per\ampere}$ (bin-averaged fit, $R^2 = 0.92$).
  \item \textbf{Theoretical $K_t\,G/r_\text{eff}$}:
        $\SI{0.632}{\newton\per\ampere}$.
  \item \textbf{Efficiency}: $\eta = 0.583/0.632 = 92\%$.
        The 8\% loss is attributed to gearbox friction, belt compliance,
        and bearing drag --- consistent with hobby-grade drivetrains.
  \item \textbf{VESC current limit}: $\sim$\SI{65}{\ampere}, producing
        a maximum wheel force of
        $F_\text{max} = 0.583 \times 65 = \SI{37.9}{\newton}$
        and a corresponding wheel torque of
        $T_\text{drive} = F_\text{max} \cdot r_\text{eff}
        = \SI{1.93}{\newton\metre}$.
\end{itemize}

\paragraph{Uncertainty.}
\begin{itemize}
  \item Motor $K_v$ from datasheet ($\pm5\%$ typical for hobby motors).
  \item $r_\text{eff}$ measurement accuracy ($\pm\SI{1}{\milli\metre}$
        $\Rightarrow \sim 2\%$ error on force).
  \item Gear efficiency not measured separately; absorbed into
        effective $K_{t,\text{eff}}$.
\end{itemize}

%%% ===================================================================
\subsection{Maximum Dynamics}\label{sec:max-dynamics}

\paragraph{Theory.}
Direct measurement of the vehicle's performance envelope:
$v_\text{max}$, $a_\text{max}$ (peak acceleration),
and $a_\text{brake}$ (peak deceleration).
These define the MPC input and state constraints.

\paragraph{Procedure.}
Command maximum speed for \SI{5}{\second} (acceleration phase),
then command zero (deceleration phase). Both ERPM-based speed and
IMU $a_x$ are recorded.

\paragraph{Data.}
\includefigure{max_dyn_speed_vs_time.pdf}{%
  Speed (odometry) vs.\ time during a maximum dynamics run.
  The braking region shows odometry dropping to zero quickly
  as the wheels lock.\label{fig:maxdyn}}

\begin{table}[H]
\centering
\caption{Maximum dynamics results.}\label{tab:maxdyn}
\begin{tabular}{lll}
\toprule
\textbf{Quantity} & \textbf{Value} & \textbf{Source} \\
\midrule
Maximum velocity $v_\text{max}$   & \SI{5.17}{\metre\per\second}           & ERPM odometry \\
Peak acceleration $a_\text{max}$  & \SI{8.01}{\metre\per\second\squared}   & IMU (smoothed) \\
Peak deceleration $a_\text{brake}$& \SI{-11.06}{\metre\per\second\squared} & IMU (smoothed) \\
\bottomrule
\end{tabular}
\end{table}

\paragraph{Validation.}
The peak acceleration can be cross-checked against the friction
coefficient.  At $a_\text{max} = \SI{8.01}{\metre\per\second\squared}$,
the implied longitudinal friction is
$\mu_x = a_\text{max}/g = 8.01/9.81 = 0.82$,
which is consistent with $\mu \approx 0.79$ from the friction test
(Section~\ref{sec:friction}).  The small difference is expected because
the IMU smoothing window captures the absolute peak, which may
briefly exceed the sustained traction limit.

The peak deceleration $|a_\text{brake}| = \SI{11.06}{\metre\per\second\squared}$
exceeds $\mu\,g \approx \SI{7.7}{\metre\per\second\squared}$, which
indicates that the IMU reading during braking includes a component from
chassis pitch (the accelerometer cannot distinguish deceleration from
forward pitch rotation under gravity).  For MPC constraints, the
conservative value $a_\text{brake} = \mu\,g \approx
\SI{7.7}{\metre\per\second\squared}$ is recommended.

\paragraph{Uncertainty.}
\begin{itemize}
  \item ERPM-derived speed is accurate during acceleration but unreliable
        during braking (wheel lock/slip).
  \item IMU-integrated velocity drifts; smoothing introduces lag.
  \item Battery voltage sag under load reduces effective $v_\text{max}$
        as the battery discharges.
\end{itemize}

%%% ===================================================================
\subsection{Steering Actuator Response}\label{sec:steering-rate}

\paragraph{Theory.}
The steering servo has a finite angular velocity. By commanding step
changes from $0 \to \delta$ and measuring the 10\%--90\% rise time
of the IMU yaw rate response:
\begin{equation}\label{eq:steer-rate}
  \dot{\delta}_\text{max} = \frac{0.8\,\delta_\text{cmd}}
                                 {\Delta t_{10\text{--}90}}.
\end{equation}
The factor 0.8 accounts for measuring the 10\% to 90\% rise region
(i.e.\ 80\% of the full step amplitude).

\paragraph{Procedure.}
Drive at constant \SI{1.5}{\metre\per\second}. Alternate steering
between $+\SI{0.3}{\radian}$ and $-\SI{0.3}{\radian}$ for 6
steps (each from the straight-ahead position),
with 1.5-second straight pauses between steps.
Record IMU yaw rate $\omega_z$ and commanded steering angle.

\paragraph{Data.}
\includefigure{steering_step_response.pdf}{%
  IMU yaw rate (solid blue, left axis) and commanded steering angle
  (dashed red, right axis) vs.\ time.
  Lines are broken at the recording gaps between
  steps.\label{fig:steer}}

\begin{table}[H]
\centering
\caption{Steering step response: 10\%--90\% rise time of the
  yaw-rate response for each transition.}\label{tab:steer}
\begin{tabular}{cS[table-format=2.1]S[table-format=3.0]S[table-format=1.2]}
\toprule
{Step} & {$\Delta\omega$ (deg/s)} & {Rise time (ms)} & {$\dot{\delta}$ (rad/s)} \\
\midrule
0 &  84.1 &  95 & 2.53 \\
1 & -79.2 &  78 & 3.08 \\
2 &  77.9 &  82 & 2.93 \\
3 & -75.9 &  76 & 3.16 \\
4 &  81.0 &  92 & 2.61 \\
5 & -75.9 &  96 & 2.50 \\
\midrule
\multicolumn{2}{c}{Average} & 87 & {$2.80^{\dagger}$} \\
\bottomrule
\end{tabular}

\smallskip\noindent${}^{\dagger}$Estimated as
$\dot{\delta}_\text{max} = 0.8 \times 0.3\;\text{rad}\;/\;
\SI{87}{\milli\second} = \SI{2.80}{\radian\per\second}$.
\end{table}

The yaw rate response shows that the combined servo and vehicle
dynamics produce consistent 10--90\% rise times of
$\sim$\SI{87}{\milli\second}.  The estimated maximum steering rate of
$\SI{2.80}{\radian\per\second}$ (auto-saved value: \SI{2.85}{\radian\per\second})
is close to the simulation default of $\SI{3.2}{\radian\per\second}$.

\paragraph{Servo gain observation.}
The VESC calibration (Section~\ref{sec:wheelbase}) shows that the
actual steering angle is approximately 86\% of the commanded angle
($\delta_\text{actual}/\delta_\text{cmd} \approx 0.86$).
This means the effective steering rate of the \emph{wheel angle} is
$0.86 \times 2.85 \approx \SI{2.45}{\radian\per\second}$.

\paragraph{Uncertainty.}
\begin{itemize}
  \item The yaw rate includes vehicle dynamics response, not only servo
        motion.  At \SI{1.5}{\metre\per\second} the dynamic lag is small
        relative to the servo speed.
  \item No direct steering angle sensor is available;
        the rate is inferred from yaw-rate dynamics.
  \item Servo speed may depend on battery voltage.
\end{itemize}

%%% ===================================================================
\subsection{Wheelbase and VESC Calibration}\label{sec:wheelbase}

\paragraph{VESC calibration.}
The VESC maps a commanded steering angle to a servo PWM position via
$s = g \cdot (-\delta_\text{cmd}) + b$.
The servo gain $g = -0.7940$, offset $b = 0.5500$, and servo limits
(\texttt{servo\_min}\,=\,0.202, \texttt{servo\_max}\,=\,0.890)
were determined with \texttt{find\_servo\_limits.py} and
\texttt{test\_steering\_calibration.py}.
The speed-to-ERPM gain was calibrated to $4550$ by comparing
VESC-reported velocity against ground-truth distance measurements.
Table~\ref{tab:vesc-cal} lists the calibrated values.

\begin{table}[H]
\centering
\caption{Calibrated VESC parameters.}\label{tab:vesc-cal}
\begin{tabular}{ll}
\toprule
\textbf{Parameter} & \textbf{Value} \\
\midrule
\texttt{speed\_to\_erpm\_gain}              & 4550 \\
\texttt{steering\_angle\_to\_servo\_gain}   & $-0.7940$ \\
\texttt{steering\_angle\_to\_servo\_offset} & $0.5500$ \\
\texttt{servo\_min} / \texttt{servo\_max}   & 0.202\,/\,0.890 \\
\bottomrule
\end{tabular}
\end{table}

With this gain, a commanded angle of
$\delta_\text{cmd} = \SI{0.3}{\radian}$ produces an actual wheel
angle of:
\begin{equation}
  \delta_\text{actual} = \Bigl|\frac{s - b}{g}\Bigr|
  = \SI{0.258}{\radian}\;(14.8°),
  \qquad
  \frac{\delta_\text{actual}}{\delta_\text{cmd}} = 0.86.
\end{equation}

\paragraph{Wheelbase test procedure.}
The effective wheelbase $L$ is measured by driving constant-radius
circles at a fixed steering angle $\delta_\text{cmd} = \SI{0.3}{\radian}$
($\delta_\text{actual} = \SI{0.258}{\radian}$)
and three speeds (1.0, 1.5, \SI{2.0}{\metre\per\second}), repeated
five times.  The circle radius $R$ is obtained from
$R = v_x / \omega_z$, where $v_x$ is the ERPM-based speed and
$\omega_z$ is the IMU yaw rate.  The wheelbase follows from the
kinematic bicycle relation:
\begin{equation}\label{eq:wheelbase}
  L = R \cdot 2\sin\!\bigl(\arctan(0.5\,\tan\delta)\bigr).
\end{equation}

\paragraph{Results.}
Table~\ref{tab:wheelbase} shows the measured effective wheelbase at
each speed.

\begin{table}[H]
\centering
\caption{Effective wheelbase $L$ (IMU-based, 5 runs mean $\pm\,\sigma$).
  CAD value: \SI{0.324}{\metre}.}\label{tab:wheelbase}
\begin{tabular}{S[table-format=1.1]
                S[table-format=1.4] @{${}\pm{}$} S[table-format=1.4]
                S[table-format=2.1]}
\toprule
{Speed (m/s)}
  & \multicolumn{2}{c}{$L_\text{measured}$ (m)}
  & {$\Delta$ from CAD (\%)} \\
\midrule
1.0 & 0.3452 & 0.0011 & +6.5 \\
1.5 & 0.3544 & 0.0013 & +9.4 \\
2.0 & 0.3690 & 0.0016 & +13.9 \\
\bottomrule
\end{tabular}
\end{table}

\includefigure{wheelbase_circle_trajectory.pdf}{%
  Wheelbase test results.
  Top-left: measured wheelbase per speed.
  Top-right: per-run consistency at \SI{1.0}{\metre\per\second}.
  Bottom-left: circle trajectories (run~5).
  Bottom-right: numerical summary.\label{fig:wheelbase}}

\paragraph{Understeer effect.}
The apparent $L$ exceeds the CAD value by
6.5\% at \SI{1.0}{\metre\per\second}, growing to 13.9\% at
\SI{2.0}{\metre\per\second}.  This speed-dependent growth is
consistent with \emph{understeer}: the front tires generate less
lateral force per unit slip angle than the rear tires, so the vehicle
tracks a larger radius than the kinematic model predicts.

The cornering stiffness test
(Section~\ref{sec:cornering}) measured
$C_{\alpha f} = \SI{49.9}{\newton\per\radian}$ and
$C_{\alpha r} = \SI{54.9}{\newton\per\radian}$ ($C_{\alpha f} <
C_{\alpha r}$), confirming the vehicle understeers.
The understeer gradient is:
\begin{equation}\label{eq:understeer}
  K_\text{us} = \frac{m\,g}{L}\Bigl(\frac{l_f}{C_{\alpha f}}
  - \frac{l_r}{C_{\alpha r}}\Bigr)
  = \SI{0.041}{\radian\per(\metre\per\second\squared)}.
\end{equation}
At \SI{1.0}{\metre\per\second} with $R \approx \SI{1.3}{\metre}$,
the lateral acceleration is
$a_y = v^2/R \approx \SI{0.76}{\metre\per\second\squared}$,
giving a steering angle loss of
$K_\text{us}\,a_y \approx 0.031\;\text{rad}\;(1.8°)$.
The kinematic formula~\eqref{eq:wheelbase} does not account for this
loss, so it over-estimates $L$.  The effect scales with $v^2$,
explaining the larger discrepancy at higher speeds.

\paragraph{Understeer-corrected wheelbase.}
If the understeer is subtracted---i.e.\ the effective steering angle
$\delta_\text{eff} = \delta - K_\text{us}\,a_y$ is used in
\eqref{eq:wheelbase}---the results at \SI{1.0}{\metre\per\second}
give $L = \SI{0.302}{\metre}$ ($-6.7\%$ from CAD).
The \emph{uncorrected} value is $\SI{0.345}{\metre}$ ($+6.5\%$),
so the CAD wheelbase of \SI{0.324}{\metre}{} lies almost exactly
between the two estimates.  This bracketing is expected: the
cornering stiffness values used for $K_\text{us}$ were measured at a
single steering angle (Section~\ref{sec:cornering}), so
the correction has significant uncertainty.  The symmetry of the
error ($\pm$6--7\%) provides confidence that the CAD value is a
good estimate of the true wheelbase.

For the MPC model, the CAD wheelbase $L = \SI{0.324}{\metre}$ is
recommended.  The understeer gradient $K_\text{us}$ is already
captured by the cornering stiffness parameters in the dynamic bicycle
model, so no additional wheelbase correction is needed.

%%% ===================================================================
\subsection{Parameter Summary}\label{sec:param-summary}

Table~\ref{tab:params} compares all measured values with the F1/10th
gym simulation defaults.  The measured values should be used to update
the MPC model for deployment on the real vehicle.

\begin{table}[H]
\centering
\caption{Identified vehicle model parameters.  --\,-- indicates the
  parameter could not be identified from the current
  data.}\label{tab:params}
\begin{tabular}{ll S[table-format=2.3] S[table-format=3.3] l l}
\toprule
{Parameter} & {Symbol} & {Measured} & {Sim default} & {Unit} & {Status} \\
\midrule
Wheelbase             & $L$
  & 0.324  & 0.324  & \si{\metre}
  & CAD\textsuperscript{d} \\
Friction coeff.\      & $\mu$
  & 0.787  & 1.049  & {--}
  & Tested \\
Front corn.\ stiff.\  & $C_{\alpha f}$
  & {--}   & 4.718  & \si{\newton\per\radian}
  & Rerun needed \\
Rear corn.\ stiff.\   & $C_{\alpha r}$
  & {--}   & 5.456  & \si{\newton\per\radian}
  & Rerun needed \\
Long.\ stiffness      & $C_x$
  & {N/A}  & 300    & \si{\newton}
  & Sensor limit\textsuperscript{a} \\
Effective $K_{t}$     & $K_{t,\text{eff}}$
  & 0.583  & {--}   & \si{\newton\per\ampere}
  & Tested \\
Max drive torque      & $T_\text{max}$
  & 1.93   & 22.9   & \si{\newton\metre}
  & Tested\textsuperscript{b} \\
Max brake decel.\     & $a_\text{brake}$
  & {7.7\textsuperscript{c}}  & 10.0   & \si{\metre\per\second\squared}
  & Tested \\
Max velocity          & $v_\text{max}$
  & 5.17   & 20.0   & \si{\metre\per\second}
  & Tested \\
Max acceleration      & $a_\text{max}$
  & 8.01   & 9.51   & \si{\metre\per\second\squared}
  & Tested \\
Max steer.\ rate      & $\dot{\delta}_\text{max}$
  & 2.85   & 3.2    & \si{\radian\per\second}
  & Tested \\
Servo gain ratio      & $\delta_\text{act}/\delta_\text{cmd}$
  & 0.86   & 1.0    & {--}
  & Calibrated\textsuperscript{e} \\
\bottomrule
\end{tabular}

\smallskip
\begin{flushleft}
\footnotesize
\textsuperscript{a}$C_x$ requires an independent body-velocity sensor
(e.g.\ lidar scan-matching) to compute the slip ratio reliably.\\
\textsuperscript{b}At the wheel ($T = F_\text{max} \times r_\text{eff}$),
at the VESC current limit of \SI{65}{\ampere}.
The simulation default is likely the motor-shaft torque
$\times\,$gear ratio, not a wheel-level measurement.\\
\textsuperscript{c}Conservative estimate: $\mu\,g$.
Raw IMU measured $\SI{11.06}{\metre\per\second\squared}$, but this
includes chassis pitch contamination.\\
\textsuperscript{d}Wheelbase test at \SI{1.0}{\metre\per\second}
yields $L = 0.345\;\text{m}$ ($+6.5\%$ from CAD);
CAD value is used (see Section~\ref{sec:wheelbase}).\\
\textsuperscript{e}Calibrated servo gain $= -0.7940$; ERPM gain $= 4550$
(see Section~\ref{sec:wheelbase}).
\end{flushleft}
\end{table}

\paragraph{Discussion.}
The measured parameters differ substantially from the simulation
defaults in several ways:
\begin{itemize}
  \item \textbf{Friction}: $\mu = 0.79$ vs.\ 1.05 in simulation.
        The real car has 25\% less grip, which will cause the MPC to
        over-estimate cornering capability if not updated.
  \item \textbf{Maximum velocity}: $v_\text{max} = \SI{5.17}{\metre\per\second}$
        vs.\ \SI{20}{\metre\per\second} in simulation.  The real vehicle's
        speed envelope is roughly 4$\times$ smaller.
  \item \textbf{Steering rate}: $\dot{\delta}_\text{max} =
        \SI{2.85}{\radian\per\second}$ vs.\ \SI{3.2}{\radian\per\second}.
        The real servo is slightly slower than assumed, but the values
        are close enough that the MPC steering rate constraint remains
        reasonable.
  \item \textbf{Motor torque} ($K_{t,\text{eff}} = \SI{0.583}{\newton\per\ampere}$):
        this value should be used instead of the theoretical
        $K_t\,G/r_\text{eff} = \SI{0.632}{\newton\per\ampere}$ to
        correctly predict wheel force from commanded motor current.
  \item \textbf{VESC calibration}: the calibrated servo gain
        ($-0.7940$) yields an actual steering angle that is 86\%
        of the commanded value.  The vehicle understeers
        ($C_{\alpha f} < C_{\alpha r}$, $K_\text{us} = 0.041$),
        which inflates the apparent turning radius and explains the
        remaining 6.5\% discrepancy in the wheelbase test
        (see Section~\ref{sec:wheelbase}).
\end{itemize}

These measurements confirm that the simulation defaults are not
representative of the physical vehicle and must be replaced before
deploying the MPC controller on the real car.

\end{document}
